% p311_1; source PDF page 311; manually verified editable table.
\begin{tikzpicture}[x=1cm,y=-1cm,inner sep=0pt,outer sep=0pt]
\path[use as bounding box] (0,0) rectangle (16.000,5.275);
\draw[line width=.4pt] (0.000,0.000) rectangle (2.700,0.815);
\node[align=center,text width=2.440cm,font={\linespread{1}\fontsize{10}{13.50}\selectfont\bfseries}] at (1.350,0.407) {总结“掐卡案”};
\draw[line width=.4pt] (2.700,0.000) rectangle (6.900,0.815);
\node[align=center,text width=3.940cm,font={\linespread{1}\fontsize{10}{13.50}\selectfont\bfseries}] at (4.800,0.407) {观点展示};
\draw[line width=.4pt] (6.900,0.000) rectangle (8.450,0.815);
\node[align=center,text width=1.290cm,font={\linespread{1}\fontsize{10}{13.50}\selectfont\bfseries}] at (7.675,0.407) {掐卡};
\draw[line width=.4pt] (8.450,0.000) rectangle (11.300,0.815);
\node[align=center,text width=2.590cm,font={\linespread{1}\fontsize{10}{13.50}\selectfont\bfseries}] at (9.875,0.407) {转账};
\draw[line width=.4pt] (11.300,0.000) rectangle (16.000,0.815);
\node[align=center,text width=4.440cm,font={\linespread{1}\fontsize{10}{13.50}\selectfont\bfseries}] at (13.650,0.407) {取款（取银行现金）};
\draw[line width=.4pt] (0.000,0.815) rectangle (2.700,3.445);
\node[align=center,text width=2.440cm,font={\linespread{1}\fontsize{10}{13.50}\selectfont}] at (1.350,2.130) {卖卡后，针对合法资金};
\draw[line width=.4pt] (2.700,0.815) rectangle (6.900,2.530);
\node[align=center,text width=3.940cm,font={\linespread{1}\fontsize{10}{13.50}\selectfont}] at (4.800,1.673) {事实占有说（多数说）};
\draw[line width=.4pt] (6.900,0.815) rectangle (8.450,2.530);
\node[align=center,text width=1.290cm,font={\linespread{1}\fontsize{10}{13.50}\selectfont}] at (7.675,1.673) {盗窃罪};
\draw[line width=.4pt] (8.450,0.815) rectangle (11.300,2.530);
\node[align=left,text width=2.590cm,font={\linespread{1}\fontsize{10}{13.50}\selectfont}] at (9.875,1.673) {掐卡定了盗窃，转账不再定罪};
\draw[line width=.4pt] (11.300,0.815) rectangle (16.000,2.530);
\node[align=left,text width=4.440cm,font={\linespread{1}\fontsize{10}{13.50}\selectfont}] at (13.650,1.673) {多数说：对银行构成犯罪。\\{}少数说：对银行不构成犯罪};
\draw[line width=.4pt] (2.700,2.530) rectangle (6.900,3.445);
\node[align=center,text width=3.940cm,font={\linespread{1}\fontsize{10}{13.50}\selectfont}] at (4.800,2.988) {法律占有说（少数说）};
\draw[line width=.4pt] (6.900,2.530) rectangle (16.000,3.445);
\node[align=center,text width=8.840cm,font={\linespread{1}\fontsize{10}{13.50}\selectfont}] at (11.450,2.988) {侵占罪};
\draw[line width=.4pt] (0.000,3.445) rectangle (2.700,5.275);
\node[align=center,text width=2.440cm,font={\linespread{1}\fontsize{10}{13.50}\selectfont}] at (1.350,4.360) {卖卡后，针对赃款};
\draw[line width=.4pt] (2.700,3.445) rectangle (6.900,4.360);
\node[align=center,text width=3.940cm,font={\linespread{1}\fontsize{10}{13.50}\selectfont}] at (4.800,3.903) {事实占有说（多数说）};
\draw[line width=.4pt] (6.900,3.445) rectangle (16.000,4.360);
\node[align=center,text width=8.840cm,font={\linespread{1}\fontsize{10}{13.50}\selectfont}] at (11.450,3.903) {盗窃罪};
\draw[line width=.4pt] (2.700,4.360) rectangle (6.900,5.275);
\node[align=center,text width=3.940cm,font={\linespread{1}\fontsize{10}{13.50}\selectfont}] at (4.800,4.818) {法律占有说（少数说）};
\draw[line width=.4pt] (6.900,4.360) rectangle (16.000,5.275);
\node[align=center,text width=8.840cm,font={\linespread{1}\fontsize{10}{13.50}\selectfont}] at (11.450,4.818) {侵占罪};
\end{tikzpicture}
